% Options for packages loaded elsewhere
\PassOptionsToPackage{unicode}{hyperref}
\PassOptionsToPackage{hyphens}{url}
\documentclass[
]{article}
\usepackage{xcolor}
\usepackage[margin=1in]{geometry}
\usepackage{amsmath,amssymb}
\setcounter{secnumdepth}{-\maxdimen} % remove section numbering
\usepackage{iftex}
\ifPDFTeX
  \usepackage[T1]{fontenc}
  \usepackage[utf8]{inputenc}
  \usepackage{textcomp} % provide euro and other symbols
\else % if luatex or xetex
  \usepackage{unicode-math} % this also loads fontspec
  \defaultfontfeatures{Scale=MatchLowercase}
  \defaultfontfeatures[\rmfamily]{Ligatures=TeX,Scale=1}
\fi
\usepackage{lmodern}
\ifPDFTeX\else
  % xetex/luatex font selection
\fi
% Use upquote if available, for straight quotes in verbatim environments
\IfFileExists{upquote.sty}{\usepackage{upquote}}{}
\IfFileExists{microtype.sty}{% use microtype if available
  \usepackage[]{microtype}
  \UseMicrotypeSet[protrusion]{basicmath} % disable protrusion for tt fonts
}{}
\makeatletter
\@ifundefined{KOMAClassName}{% if non-KOMA class
  \IfFileExists{parskip.sty}{%
    \usepackage{parskip}
  }{% else
    \setlength{\parindent}{0pt}
    \setlength{\parskip}{6pt plus 2pt minus 1pt}}
}{% if KOMA class
  \KOMAoptions{parskip=half}}
\makeatother
\usepackage{longtable,booktabs,array}
\usepackage{calc} % for calculating minipage widths
% Correct order of tables after \paragraph or \subparagraph
\usepackage{etoolbox}
\makeatletter
\patchcmd\longtable{\par}{\if@noskipsec\mbox{}\fi\par}{}{}
\makeatother
% Allow footnotes in longtable head/foot
\IfFileExists{footnotehyper.sty}{\usepackage{footnotehyper}}{\usepackage{footnote}}
\makesavenoteenv{longtable}
\setlength{\emergencystretch}{3em} % prevent overfull lines
\providecommand{\tightlist}{%
  \setlength{\itemsep}{0pt}\setlength{\parskip}{0pt}}
\usepackage{bookmark}
\IfFileExists{xurl.sty}{\usepackage{xurl}}{} % add URL line breaks if available
\urlstyle{same}
\hypersetup{
  hidelinks,
  pdfcreator={LaTeX via pandoc}}

\author{}
\date{}

\begin{document}

\section{Selectors for the Born Pairing Among Simultaneously Unitarily
Invariant Bilinear
Pairings}\label{selectors-for-the-born-pairing-among-simultaneously-unitarily-invariant-bilinear-pairings}

\textbf{Author:} Paulina Janowska \textbf{Affiliation:} Independent
researcher, Poznan, Poland \textbf{Date:} 2026-07-24 \textbf{Version:}
4.7 (distinguishability selector)

\begin{center}\rule{0.5\linewidth}{0.5pt}\end{center}

\subsection{Abstract}\label{abstract}

Let \(\widetilde B:\mathrm{Herm}_d\times\mathrm{Herm}_d\to\mathbb R\) be
a real-bilinear form, invariant under simultaneous unitary conjugation,
\(\widetilde B(UEU^\dagger,U\rho U^\dagger)=\widetilde B(E,\rho)\),
whose restriction to effects and normalized states
\(B:=\widetilde B|_{\mathrm{Eff}_d\times\mathcal D_d}\) takes values in
\([0,1]\) and is normalized, \(B(I,\rho)=1\) for
\(\rho\in\mathcal D_d\). The classification is elementary: by the
decomposition \(\mathrm{Herm}_d=\mathbb RI\oplus\mathrm{Herm}_d^0\) into
trivial and adjoint irreducible components, every such form is
\[\widetilde B(E,\rho)=a\,\mathrm{Tr}(E\rho)+b\,\mathrm{Tr}(E)\,\mathrm{Tr}(\rho)\]
with real \(a,b\) and \(a+bd=1\); positivity forces \(a\in[-1/(d-1),1]\)
(for \(d\geq 2\)). The Born rule \(B(E,\rho)=\mathrm{Tr}(E\rho)\)
corresponds to \(a=1,b=0\), but is \emph{not} selected by invariance,
bilinearity, positivity, and normalization alone. We give five selector
conditions that do select Born, each after the classification:

\begin{itemize}
\tightlist
\item
  \textbf{Theorem 1 (one impossible event).} If \(B\) vanishes on a
  single pair of orthogonal rank-one projectors \(P,Q\)
  (i.e.~\(B(Q,P)=0\) with \(PQ=0\)), then \(b=0\) and
  \(B=\mathrm{Tr}(E\rho)\).
\item
  \textbf{Theorem 2 (rank-one product monoidality).} If the family
  \(\{B_d\}\) is multiplicative on rank-one product effects and
  normalized product states, and non-trivial (\(a_d\neq 0\)), then
  \(B_d=\mathrm{Tr}(E\rho)\) for every \(d\geq 2\).
\item
  \textbf{Theorem 3 (pure-state self-certainty).} If \(B(\psi,\psi)=1\)
  for every pure state \(\psi\) (viewed as a rank-one projector, hence
  both a state and an effect), then \(a=1\) and
  \(B=\mathrm{Tr}(E\rho)\).
\item
  \textbf{Theorem 4 (dimension-independence of pure transitions).} If
  the transition probability between two pure states depends only on the
  two states and not on the ambient dimension \(d\) into which they are
  embedded, then \(a_d=1\) for all \(d\geq 2\).
\item
  \textbf{Theorem 5 (perfect distinguishability).} If two orthogonal
  pure states are perfectly distinguishable by some effect (\(B(E,P)=1\)
  and \(B(E,Q)=0\)), then for \(d\geq 3\) necessarily \(a=1\) (Born);
  for \(d=2\) one has \(a\in\{1,-1\}\), and the anti-Born endpoint
  \(a=-1\) is killed by any of Theorems 1--3.
\end{itemize}

The selectors are elementary: each is a one-line consequence of the
classification plus a single operational test. Theorems 1 and 3 are
equivalent under the classification (Proposition 3). Theorems 1--3
appear to be unrecorded as Born selectors \emph{within the invariant
bilinear class}; Theorem 4 is a restricted restatement of Hossenfelder's
dimension-independence argument {[}22{]}. We claim them as corollaries,
not as deep new theorems. We do not derive the Hermitian structure or
quantum theory; Nakahira {[}8{]} already derives the full complex
quantum package (including Born) from two OPT postulates, so the present
selectors live one layer below that reconstruction.

\begin{center}\rule{0.5\linewidth}{0.5pt}\end{center}

\subsection{1. Introduction}\label{introduction}

\subsubsection{1.1 The question}\label{the-question}

The Born rule \(p(E|\rho)=\mathrm{Tr}(E\rho)\) is the standard pairing
of a quantum state \(\rho\) with an effect \(E\). It is \emph{not} the
unique bilinear, positive, unitarily invariant, normalized pairing: the
one-parameter family
\(B_a(E,\rho)=a\,\mathrm{Tr}(E\rho)+(1-a)/d\cdot\mathrm{Tr}(E)\mathrm{Tr}(\rho)\)
with \(a\in[-1/(d-1),1]\) satisfies all four conditions, and for
\(a\neq 1\) is not Born. This family is the full classification (Lemma 2
below). The question is: which simple additional condition selects Born?

\subsubsection{1.2 What this paper is}\label{what-this-paper-is}

This paper is:

\begin{itemize}
\tightlist
\item
  a self-contained classification of simultaneously unitarily invariant
  real-bilinear forms on \(\mathrm{Herm}_d\) (Lemma 2), via the
  decomposition into trivial and adjoint irreducible components;
\item
  \textbf{five selector theorems (Theorems 1--5)} that select the Born
  pairing from the classified family under minimal operational
  conditions;
\item
  Proposition 3: local equivalence of Theorems 1 and 3;
\item
  a small classification lemma for the Gram map on complex amplitudes
  (Lemma 1), included as scaffolding;
\item
  an honest statement of the circularity boundary and of the relation to
  Galley--Masanes {[}9{]} and Nakahira {[}8{]}.
\end{itemize}

This paper is not:

\begin{itemize}
\tightlist
\item
  a derivation of the Born rule from operational principles;
\item
  a derivation of the Hermitian structure, the effect cone, or the
  complex amplitude representation;
\item
  a claim that the selectors are deep new theorems (they are elementary
  corollaries of the classification);
\item
  a reconstruction of quantum theory.
\end{itemize}

\subsubsection{1.3 Related prior work}\label{related-prior-work}

The decomposition \(\mathrm{Herm}_d=\mathbb RI\oplus\mathrm{Herm}_d^0\)
and the uniqueness of invariant forms on irreducible representations are
standard representation theory. Galley and Masanes {[}9{]} classify
\emph{all} alternatives to the measurement postulates that keep pure
states as rays in \(\mathbb{CP}^{d-1}\) and unitary dynamics, subject to
finite-dimensional mixed-state spaces; their alternatives correspond to
representations of the unitary group and are far broader than the
bilinear class studied here. In particular, unrestricted non-Born
alternatives in their catalogue violate bit symmetry. Our Lemma 2 is the
thin slice of that landscape in which outcome probabilities extend to a
real-bilinear form on \(\mathrm{Herm}_d\times\mathrm{Herm}_d\): only the
family \(B_{a,b}\) survives, and Theorems 1--4 then select Born inside
that slice. Hossenfelder {[}22{]} uses dimension-independence of
transition probabilities on the complex sphere (a pure-state continuous
setting that already includes orthogonal vanishing). Lax {[}23{]} uses
Fisher-non-expansion and Cramér-Rao bounds. Gleason {[}14{]} uses
additivity over orthogonal decompositions; Busch {[}15{]} extends to
POVMs and the qubit case. Chiribella, D'Ariano, and Perinotti {[}4{]}
derive quantum theory from purification; Barnum and Wilce {[}5{]} from
local tomography plus Jordan structure; Tull {[}6{]} and Selby,
Scandolo, and Coecke {[}7{]} categorically; Nakahira {[}8{]} from two
operational principles (local equivalence + ES purification), recovering
the full complex package including Born --- so our selectors are not a
competing OPT reconstruction. Procesi {[}2{]} gives the tensor first
fundamental theorem for polynomial equivariant maps (our Lemma 1 is a
continuous special case, not formally subsumed). Massart and Absil
{[}3{]} treat the quotient geometry of fixed-rank PSD matrices.

\subsubsection{1.4 Outline}\label{outline}

Section 2 fixes notation. Section 3 states and proves Lemma 1 (Gram map
classification) as scaffolding. Section 4 records the fibre and quotient
geometry. Section 5 states the classification of invariant bilinear
pairings (Lemma 2) with a self-contained proof via the irreducible
decomposition. Section 6 gives the five selector theorems. Section 7
discusses limits, circularity, and the status of the selectors. Section
8 confronts an open conjecture with Nakahira's result. Section 9 records
an informal operational extension.

\begin{center}\rule{0.5\linewidth}{0.5pt}\end{center}

\subsection{2. Preliminaries}\label{preliminaries}

Fix \(d\geq 2\) (the case \(d=1\) is trivial and excluded from the
selectors). Write \(\mathrm{Herm}_d\) for Hermitian \(d\times d\)
matrices, \(\mathrm{PSD}_d\) for the positive semidefinite cone,
\(\mathcal D_d=\{\rho\succeq 0:\mathrm{Tr}\,\rho=1\}\) for density
matrices, and \(\mathrm{Eff}_d=\{E:0\leq E\leq I\}\) for effects. The
unitary group \(U(d)\) acts on \(\mathrm{Herm}_d\) by simultaneous
conjugation \(E\mapsto UEU^\dagger\).

A key distinction: for a scalar-valued form \(B(E,\rho)\), the condition
\(B(UEU^\dagger,U\rho U^\dagger)=B(E,\rho)\) is \textbf{invariance} (the
value is unchanged), not covariance (covariance applies to
operator-valued maps whose output transforms). We use ``invariant''
throughout.

\begin{center}\rule{0.5\linewidth}{0.5pt}\end{center}

\subsection{3. The Gram map
(scaffolding)}\label{the-gram-map-scaffolding}

\subsubsection{3.1 Statement}\label{statement}

\textbf{Lemma 1 (Gram map uniqueness).} Let
\(q:M_{d\times r}(\mathbb C)\to\mathrm{Herm}_d\) be a continuous map
satisfying, for every \(W\), \(M\in GL(d,\mathbb C)\), \(U\in U(r)\):

\begin{enumerate}
\def\labelenumi{\arabic{enumi}.}
\tightlist
\item
  \(q(WU)=q(W)\) (right gauge invariance);
\item
  \(q(MW)=M\,q(W)\,M^\dagger\) (left congruence covariance);
\item
  \(q(W)\succeq 0\) (positivity).
\end{enumerate}

Then \(q(W)=c\,WW^\dagger\) for some \(c\geq 0\).

\subsubsection{3.2 Proof}\label{proof}

\textbf{Step 1 (rank reduction).} Let \(k=\mathrm{rank}\,W\) and
\(J_k=\begin{pmatrix}I_k&0\\0&0\end{pmatrix}\). Every rank-\(k\) matrix
has a decomposition \(W=MJ_kU\) with \(M\in GL(d,\mathbb C)\),
\(U\in U(r)\) (full SVD, singular values absorbed into \(M\)). By (1)
then (2), \(q(W)=M\,q(J_k)\,M^\dagger\), so \(q\) on the rank-\(k\)
stratum is determined by \(q(J_k)\).

\textbf{Step 2 (stabilizer).} For
\(S=\begin{pmatrix}I_k&B\\0&C\end{pmatrix}\) with
\(B\in M_{k\times(d-k)}\), \(C\in GL(d-k)\), we have \(SJ_k=J_k\). By
(2), \(q(J_k)=Sq(J_k)S^\dagger\). Writing
\(q(J_k)=\begin{pmatrix}X&Y\\Y^\dagger&Z\end{pmatrix}\), the constraint
for all \(B,C\) forces \(Z=0\) (from \(CZC^\dagger=Z\)) and then \(Y=0\)
(from \((Y+BZ)C^\dagger=Y\)). So
\(q(J_k)=\begin{pmatrix}A_k&0\\0&0\end{pmatrix}\), \(A_k\succeq 0\).

\textbf{Step 3 (Schur).} For \(V\in U(k)\),
\(L_V=\mathrm{diag}(V,I_{d-k})\) and
\(R_V=\mathrm{diag}(V^\dagger,I_{r-k})\) satisfy \(L_VJ_kR_V=J_k\). By
(1),(2), \(VA_kV^\dagger=A_k\) for all \(V\in U(k)\). Since \(U(k)\)
acts irreducibly on \(\mathbb C^k\), Schur's lemma gives \(A_k=c_kI_k\).

\textbf{Step 4 (continuity).}
\(W_\varepsilon=\mathrm{diag}(s_1,\ldots,s_k,\varepsilon,0,\ldots)\to W_0=\mathrm{diag}(s_1,\ldots,s_k,0,\ldots)\)
as \(\varepsilon\to 0\). Continuity of \(q\) gives
\(c_{k+1}W_0W_0^\dagger=q(W_\varepsilon)\to q(W_0)=c_kW_0W_0^\dagger\),
so \(c_{k+1}=c_k\). Thus \(q(W)=c\,WW^\dagger\). \(\blacksquare\)

\subsubsection{3.3 Why unitary covariance is too
weak}\label{why-unitary-covariance-is-too-weak}

\(q_n(W)=(WW^\dagger)^n\) is unitarily covariant for all \(n\geq 1\)
(conjugation commutes with powers) but fails full congruence covariance
for \(n>1\). Congruence under \(GL(d,\mathbb C)\), not just \(U(d)\), is
what selects degree two.

\begin{center}\rule{0.5\linewidth}{0.5pt}\end{center}

\subsection{4. Fibres and quotient
geometry}\label{fibres-and-quotient-geometry}

\textbf{Proposition 1.} For \(W,V\in M_{d\times r}(\mathbb C)\) with
\(WW^\dagger=VV^\dagger\), there exists \(U\in U(r)\) with \(V=WU\).

\emph{Proof.} Write full singular value decompositions
\(W=L_W\Sigma_W R_W^\dagger\), \(V=L_V\Sigma_V R_V^\dagger\) with
\(L_W,L_V\in U(d)\), \(R_W,R_V\in U(r)\), \(\Sigma\) rectangular
diagonal. The equality \(WW^\dagger=VV^\dagger\) is
\(L_W\Sigma_W\Sigma_W^\dagger L_W^\dagger=L_V\Sigma_V\Sigma_V^\dagger L_V^\dagger\).
This is an equality of spectral decompositions of the same PSD matrix
\(\rho:=WW^\dagger=VV^\dagger\). By uniqueness of the spectral
decomposition:

\begin{itemize}
\tightlist
\item
  the eigenvalue multisets coincide, so the non-zero singular values of
  \(W\) and \(V\) coincide, and after reordering
  \(\Sigma_W=\Sigma_V=:\Sigma\);
\item
  the left unitary factors \(L_W,L_V\) are equal up to block-diagonal
  unitary on the degenerate eigenspaces of \(\rho\) (on each eigenspace
  of \(\rho\) with a given eigenvalue, the choice of eigenbasis is free
  up to a unitary). Absorb this block-diagonal freedom into \(R_W\)
  (redefining \(R_W\) by left-multiplication with the block-diagonal
  unitary), so that \(L_W=L_V=:L\).
\end{itemize}

With \(W=L\Sigma R_W^\dagger\) and \(V=L\Sigma R_V^\dagger\), we have
\(V=W(R_WR_V^\dagger)\) and \(U:=R_WR_V^\dagger\in U(r)\). \(\square\)

Thus
\(M_{d\times r}(\mathbb C)/U(r)\cong\{\rho\succeq 0:\mathrm{rank}\,\rho\leq r\}\)
(all ranks), standard quotient geometry {[}3{]}.

\textbf{Proposition 2 (monoidality after normalization).} With
\(\mathrm{Tr}\,q(W)=\mathrm{Tr}(WW^\dagger)\) forcing \(c=1\),
\(q(W)=WW^\dagger\) is strictly monoidal:
\(q(W_A\otimes W_B)=q(W_A)\otimes q(W_B)\).

\begin{center}\rule{0.5\linewidth}{0.5pt}\end{center}

\subsection{5. Classification of invariant bilinear
pairings}\label{classification-of-invariant-bilinear-pairings}

\subsubsection{5.1 Statement}\label{statement-1}

\textbf{Lemma 2.} Let
\(\widetilde B:\mathrm{Herm}_d\times\mathrm{Herm}_d\to\mathbb R\) be
real-bilinear and invariant under simultaneous unitary conjugation:
\(\widetilde B(UEU^\dagger,U\rho U^\dagger)=\widetilde B(E,\rho)\) for
all \(U\in U(d)\). Then there exist real \(a,b\) (depending on \(d\))
such that
\[\widetilde B(E,\rho)=a\,\mathrm{Tr}(E\rho)+b\,\mathrm{Tr}(E)\,\mathrm{Tr}(\rho).\]

\subsubsection{5.2 Self-contained proof via irreducible
decomposition}\label{self-contained-proof-via-irreducible-decomposition}

Decompose \(\mathrm{Herm}_d=\mathbb RI\oplus\mathrm{Herm}_d^0\), where
\(\mathbb RI=\{\lambda I\}\) and
\(\mathrm{Herm}_d^0=\{H\in\mathrm{Herm}_d:\mathrm{Tr}\,H=0\}\). Under
\(U(d)\) conjugation, \(\mathbb RI\) carries the trivial representation
and \(\mathrm{Herm}_d^0\) carries the adjoint representation.

\textbf{Claim 1 (real-irreducibility of \(\mathrm{Herm}_d^0\)).} The
adjoint representation of \(U(d)\) on \(\mathrm{Herm}_d^0\) is
real-irreducible.

\emph{Proof of Claim 1.} Complexify:
\(\mathrm{Herm}_d^0\otimes\mathbb C=\mathfrak{sl}_d(\mathbb C)\), which
is the complex irreducible adjoint representation of \(U(d)\)
(equivalently of \(SU(d)\), since the central \(U(1)\) acts trivially on
traceless matrices). A complex-irreducible representation whose real
form is reducible must be of complex type (admit a \(U(d)\)-invariant
complex structure); \(\mathfrak{sl}_d\) is of real type for \(d\geq 2\)
(it carries the invariant real form \(\mathrm{Herm}_d^0\) and no
invariant complex structure), so \(\mathrm{Herm}_d^0\) is
real-irreducible. \(\square\)

\textbf{Claim 2 (no cross-terms).} Any invariant bilinear form
\(\widetilde B\) satisfies \(\widetilde B(S,T_0)=\widetilde B(T_0,S)=0\)
for \(S\in\mathbb RI\), \(T_0\in\mathrm{Herm}_d^0\).

\emph{Proof of Claim 2.} Fix \(S=\lambda I\in\mathbb RI\). The map
\(\phi_\lambda:\mathrm{Herm}_d^0\to\mathbb R\),
\(\phi_\lambda(T_0)=\widetilde B(\lambda I,T_0)\), is
\(U(d)\)-equivariant (by invariance of \(\widetilde B\):
\(\widetilde B(U\lambda IU^\dagger,UT_0U^\dagger)=\widetilde B(\lambda I,T_0)\),
and \(U\lambda IU^\dagger=\lambda I\)). Since \(\mathbb RI\) carries the
trivial representation, \(\phi_\lambda\) is an intertwiner from the
adjoint (non-trivial irreducible) to the trivial representation. By
Schur's lemma (real form: any equivariant linear map between
non-isomorphic irreducible real representations is zero),
\(\phi_\lambda=0\), giving \(\widetilde B(S,T_0)=0\). The reversed
cross-term \(\widetilde B(T_0,S)\) vanishes by the same argument with
the two slots interchanged (or by symmetry of the invariant form, if
\(B\) is symmetric; in general, apply the argument to
\(\widetilde B'(E,\rho):=\widetilde B(\rho,E)\), which is also
invariant). \(\square\)

\textbf{Claim 3 (uniqueness on each component).} The space of invariant
bilinear forms on \(\mathbb RI\times\mathbb RI\) is one-dimensional, and
on \(\mathrm{Herm}_d^0\times\mathrm{Herm}_d^0\) is one-dimensional.

\emph{Proof of Claim 3.} On \(\mathbb RI\times\mathbb RI\): identify
\(\mathbb RI\cong\mathbb R\) via \(\lambda I\mapsto\lambda\); the
trivial action leaves any bilinear form
\(\lambda\mu\mapsto c\lambda\mu\) invariant, a one-dimensional space.

On \(\mathrm{Herm}_d^0\times\mathrm{Herm}_d^0\): let
\(V:=\mathrm{Herm}_d^0\). The adjoint representation of \(U(d)\) on
\(V\) is absolutely irreducible (complexification \(\mathfrak{sl}_d\) is
complex-irreducible), so by the real Schur lemma its commutant is
\(\mathrm{End}_{U(d)}(V)\cong\mathbb R\) (scalar endomorphisms only).
The trace form \(\langle A,B\rangle:=\mathrm{Tr}(AB)\) is a
non-degenerate \(U(d)\)-invariant bilinear form on \(V\) (invariance by
cyclicity: \(\mathrm{Tr}(UAU^\dagger\,UBU^\dagger)=\mathrm{Tr}(AB)\)),
hence gives a \(U(d)\)-equivariant isomorphism \(\psi:V\to V^*\),
\(\psi(A)=\langle A,\cdot\rangle\). Now any invariant bilinear form
\(B_0\) on \(V\) corresponds, via \(\psi^{-1}\) on the second slot, to a
\(U(d)\)-equivariant endomorphism \(\Phi:V\to V\) defined by
\(B_0(A,B)=\langle\Phi(A),B\rangle=\mathrm{Tr}(\Phi(A)B)\).
(Equivariance of \(\Phi\) follows from invariance of \(B_0\) and
\(\langle\cdot,\cdot\rangle\).) Since
\(\mathrm{End}_{U(d)}(V)=\mathbb R\), \(\Phi=\alpha\,\mathrm{id}\) for
some \(\alpha\in\mathbb R\), giving
\(B_0(A,B)=\alpha\,\mathrm{Tr}(AB)\). Thus the space of invariant
bilinear forms on \(V\times V\) is one-dimensional, spanned by the trace
form. \(\square\)

\textbf{Remark.} The key facts are: (i) the adjoint representation on
\(\mathrm{Herm}_d^0\) is absolutely irreducible, so
\(\mathrm{End}_{U(d)}(V)\cong\mathbb R\); (ii) the trace form is a
non-degenerate invariant bilinear form giving \(V\cong V^*\). Together
these force the invariant bilinear form to be a scalar multiple of the
trace. These are standard results in invariant theory for compact
groups. \(\square\)

Combining Claims 1-3:
\(\widetilde B(E,\rho)=b\,(\mathrm{Tr}(E)/d)(\mathrm{Tr}(\rho)/d)\cdot d^2 + a\,\mathrm{Tr}(E_0\rho_0)\)
where \(E_0,\rho_0\) are traceless parts; using
\(\mathrm{Tr}(E\rho)=\mathrm{Tr}(E_0\rho_0)+\mathrm{Tr}(E)\mathrm{Tr}(\rho)/d\)
and absorbing the constant into the definition of \(b\), we obtain
\(\widetilde B(E,\rho)=a\,\mathrm{Tr}(E\rho)+b\,\mathrm{Tr}(E)\mathrm{Tr}(\rho)\).
\(\blacksquare\)

\subsubsection{5.3 Positivity and
normalization}\label{positivity-and-normalization}

Restrict to \(B:=\widetilde B|_{\mathrm{Eff}_d\times\mathcal D_d}\).
Normalization \(B(I,\rho)=1\) for \(\mathrm{Tr}\,\rho=1\) gives
\(a+bd=1\), so \(b=(1-a)/d\). Positivity on all effects and normalized
states forces \(a\in[-1/(d-1),1]\): for \(a<0\), the minimum of
\(a\,\mathrm{Tr}(E\rho)\) over normalized \(\rho\) is
\(a\,\lambda_{\max}(E)\), and combining with \(b\,\mathrm{Tr}(E)\) gives
minimum \((a+b)\mathrm{Tr}(E)\) on rank-one \(E\), requiring
\(a+b\geq 0\), i.e.~\(a\geq -1/(d-1)\); for \(a\geq 0\), the minimum is
\(b\,\mathrm{Tr}(E)\geq 0\), requiring \(a\leq 1\).

The Born rule is \(a=1\) (so \(b=0\)). It is \emph{not} selected by
invariance, bilinearity, positivity, and normalization alone.

\begin{center}\rule{0.5\linewidth}{0.5pt}\end{center}

\subsection{6. Selector theorems}\label{selector-theorems}

\subsubsection{6.1 Theorem 1: one impossible
event}\label{theorem-1-one-impossible-event}

\textbf{Theorem 1.} Let \(d\geq 2\) and let
\(\widetilde B:\mathrm{Herm}_d\times\mathrm{Herm}_d\to\mathbb R\) be
real-bilinear and unitarily invariant, with
\(B:=\widetilde B|_{\mathrm{Eff}_d\times\mathcal D_d}\) normalized
(\(B(I,\rho)=1\) for \(\rho\in\mathcal D_d\)). Suppose there exist two
orthogonal rank-one projectors \(P,Q\) (i.e.~\(PQ=0\)) such that
\(B(Q,P)=0\). Then \(B(E,\rho)=\mathrm{Tr}(E\rho)\) on
\(\mathrm{Eff}_d\times\mathcal D_d\).

\emph{Proof.} By Lemma 2,
\(B(E,\rho)=a\,\mathrm{Tr}(E\rho)+b\,\mathrm{Tr}(E)\mathrm{Tr}(\rho)\)
with \(a+bd=1\). For rank-one projectors \(P,Q\) with \(PQ=0\):
\(\mathrm{Tr}(QP)=0\), \(\mathrm{Tr}(Q)=\mathrm{Tr}(P)=1\). So
\(B(Q,P)=a\cdot 0+b\cdot 1\cdot 1=b\). The hypothesis \(B(Q,P)=0\)
forces \(b=0\), hence \(a=1\), hence \(B(E,\rho)=\mathrm{Tr}(E\rho)\).
\(\blacksquare\)

\textbf{Remark.} The hypothesis is stated for a single pair \(P,Q\). By
unitary invariance, \(B\) vanishes on every pair of orthogonal rank-one
projectors (any such pair is \(UPU^\dagger,UQU^\dagger\) for some
\(U\)). So one calibrated impossible event --- ``this measurement
returns zero on a state orthogonal to it'' --- selects Born within the
invariant bilinear class.

\subsubsection{6.2 Theorem 2: rank-one product
monoidality}\label{theorem-2-rank-one-product-monoidality}

\textbf{Definition 1 (rank-one product monoidality).} A family
\(\{B_d\}_{d\geq 2}\) is \emph{rank-one product monoidal} if for every
\(d_A,d_B\geq 2\), every normalized \(\rho_A\in\mathcal D_{d_A}\),
\(\rho_B\in\mathcal D_{d_B}\), and rank-one projectors
\(P_A=|e\rangle\langle e|\), \(P_B=|f\rangle\langle f|\),
\[B_{d_Ad_B}(P_A\otimes P_B,\;\rho_A\otimes\rho_B) = B_{d_A}(P_A,\rho_A)\,B_{d_B}(P_B,\rho_B).\]

\textbf{Definition 2 (non-triviality).} \(B_d\) is non-trivial if
\(a_d\neq 0\).

\textbf{Theorem 2.} Let \(\{B_d\}_{d\geq 2}\) be a family of normalized
restrictions of unitarily invariant real-bilinear forms on
\(\mathrm{Herm}_d\). If the family is rank-one product monoidal and
non-trivial on every system, then \(b_d=0\) and \(a_d=1\) for every
\(d\geq 2\), i.e.~\(B_d(E,\rho)=\mathrm{Tr}(E\rho)\).

\emph{Proof.} Fix \(d_A,d_B\geq 2\). Set
\(X=\langle e|\rho_A|e\rangle\in[0,1]\),
\(Y=\langle f|\rho_B|f\rangle\in[0,1]\) (these vary independently as
\(\rho_A,\rho_B\) range over \(\mathcal D\)). For rank-one projectors,
\(\mathrm{Tr}(P_A\rho_A)=X\), \(\mathrm{Tr}(P_B\rho_B)=Y\),
\(\mathrm{Tr}(P_A)=\mathrm{Tr}(P_B)=1\),
\(\mathrm{Tr}(\rho_A)=\mathrm{Tr}(\rho_B)=1\). The monoidality identity
becomes
\[a_{d_Ad_B}XY + b_{d_Ad_B}\cdot 1\cdot 1 = (a_{d_A}X+b_{d_A})(a_{d_B}Y+b_{d_B}).\]
Expanding the right side:
\(a_{d_A}a_{d_B}XY+a_{d_A}b_{d_B}X+b_{d_A}a_{d_B}Y+b_{d_A}b_{d_B}\).
Matching coefficients of the polynomial in \(X,Y\in[0,1]\) (the identity
holds for all such \(X,Y\), so coefficients match): - \(XY\):
\(a_{d_Ad_B}=a_{d_A}a_{d_B}\); - \(X\): \(0=a_{d_A}b_{d_B}\); - \(Y\):
\(0=b_{d_A}a_{d_B}\); - const: \(b_{d_Ad_B}=b_{d_A}b_{d_B}\).

Non-triviality (\(a_{d_A},a_{d_B}\neq 0\)) forces \(b_{d_A}=b_{d_B}=0\).
Then normalization gives \(a_{d_A}=a_{d_B}=1\). \(\blacksquare\)

\textbf{Remark (degenerate case).} Without non-triviality,
\(B_d(E,\rho)=\mathrm{Tr}(E)\mathrm{Tr}(\rho)/d\) (\(a_d=0,b_d=1/d\)) is
rank-one product monoidal. Non-triviality excludes this ``classical''
pairing that ignores state-effect alignment.

\subsubsection{6.3 Theorem 3: pure-state
self-certainty}\label{theorem-3-pure-state-self-certainty}

\textbf{Theorem 3.} Let \(d\geq 2\) and let \(B\) be as in Theorem 1
(normalized restriction of a unitarily invariant real-bilinear form).
Suppose that for every pure state \(\psi\in\mathcal D_d\) (equivalently,
every rank-one projector \(P\)), \[B(P,P)=1.\] Then
\(B(E,\rho)=\mathrm{Tr}(E\rho)\).

\emph{Proof.} A pure state is a rank-one projector \(P\), which is both
an effect and a normalized state. By Lemma 2,
\(B(P,P)=a\,\mathrm{Tr}(P^2)+b\,\mathrm{Tr}(P)^2=a\cdot 1+b\cdot 1=a+b\).
The hypothesis \(B(P,P)=1\) gives \(a+b=1\). Combined with normalization
\(a+bd=1\):
\[a+b=a+\frac{1-a}{d}=\frac{a(d-1)+1}{d}=1\quad\Rightarrow\quad a(d-1)+1=d\quad\Rightarrow\quad a(d-1)=d-1.\]
Since \(d\geq 2\), \(a=1\), hence \(b=0\). \(\blacksquare\)

\textbf{Remark.} Operationally: a pure state is certain about itself.
Measuring the projector \(P\) on the state \(P\) must return ``yes''
with probability 1. This single numerical condition on the diagonal of
pure states kills the entire \(b\)-family. Note that \(P\) must be
admissible both as effect and as state; this is automatic for rank-one
projectors in the standard setup.

\subsubsection{6.4 Theorem 4: dimension-independence of pure
transitions}\label{theorem-4-dimension-independence-of-pure-transitions}

\textbf{Definition 3 (dimension-independence).} A family
\(\{B_d\}_{d\geq 2}\) is \emph{dimension-independent on pure
transitions} if, whenever two pure states \(\psi,\phi\) of a
\(d_0\)-dimensional system are embedded isometrically into
\(\mathbb C^d\) for any \(d\geq d_0\) (as \(\iota(\psi),\iota(\phi)\)),
the transition probability is independent of \(d\):
\[B_d(\iota(\phi),\iota(\psi))=B_{d_0}(\phi,\psi).\]

\textbf{Theorem 4.} Let \(\{B_d\}_{d\geq 2}\) be a family of normalized
restrictions of unitarily invariant real-bilinear forms. If the family
is dimension-independent on pure transitions, then \(a_d=1\) and
\(B_d=\mathrm{Tr}(E\rho)\) for every \(d\geq 2\).

\emph{Proof.} Let \(P,Q\) be rank-one projectors in \(\mathbb C^2\) with
\(\mathrm{Tr}(PQ)=x\in[0,1]\). Embed them into \(\mathbb C^d\) for any
\(d\geq 2\). Then
\[B_d(Q,P)=a_d\,x+b_d\cdot 1\cdot 1=a_d\,x+\frac{1-a_d}{d}.\]
Dimension-independence requires that \(a_d\,x+(1-a_d)/d\) is independent
of \(d\) for all \(x\in[0,1]\) and all \(d\geq 2\). Taking two values of
\(x\) (e.g.~\(x=0\) and \(x=1\)): - \(x=0\): \((1-a_d)/d = c_0\)
(constant in \(d\)); - \(x=1\): \(a_d+(1-a_d)/d = c_1\) (constant in
\(d\)).

From \(x=0\): \(1-a_d=c_0\,d\) for all \(d\geq 2\). The only way a
linear function of \(d\) can hold for all \(d\geq 2\) with
\(a_d\in[-1/(d-1),1]\) is \(c_0=0\), hence \(a_d=1\) for all \(d\). (If
\(c_0\neq 0\), then \(a_d=1-c_0 d\) eventually leaves the positivity
interval \([-1/(d-1),1]\).) \(\blacksquare\)

\textbf{Remark.} This is a restricted restatement of Hossenfelder's
dimension-independence argument {[}22{]}, placed \emph{after} the
invariant bilinear classification. Hossenfelder works directly with
continuous unitarily invariant transition probabilities \(P_N\) on the
complex unit sphere, assumes \(N\)-independence and basis consistency
(which already includes \(P=0\) on orthogonal pairs --- our Theorem 1
--- and \(\sum_i P(\Psi\to e_i)=1\)), and concludes
\(P=|\langle\Psi|\Phi\rangle|^2\). Our Theorem 4 isolates only the
dimension-independence step inside the coarser bilinear class
\(B_{a,b}\), where the orthogonal-vanishing step is separated as Theorem
1. We include Theorem 4 for catalogue completeness and to make the
relation to {[}22{]} explicit, not as an independent claim of priority.

\subsubsection{6.5 Theorem 5: perfect
distinguishability}\label{theorem-5-perfect-distinguishability}

\textbf{Definition 4.} Two pure states \(P,Q\) are \emph{perfectly
distinguishable} by an effect \(E\in\mathrm{Eff}_d\) if \(B(E,P)=1\) and
\(B(E,Q)=0\).

\textbf{Theorem 5.} Let \(d\geq 2\) and let \(B\) be a normalized
restriction of a unitarily invariant real-bilinear form as in Lemma 2,
with \(a\in[-1/(d-1),1]\). Suppose there exist orthogonal pure states
\(P,Q\) and an effect \(E\) such that \(B(E,P)=1\) and \(B(E,Q)=0\).
Then:

\begin{enumerate}
\def\labelenumi{\arabic{enumi}.}
\tightlist
\item
  if \(d\geq 3\), necessarily \(a=1\) (hence \(B\) is Born);
\item
  if \(d=2\), necessarily \(a\in\{1,-1\}\); the anti-Born endpoint
  \(a=-1\) (so \(b=1\)) satisfies distinguishability via \(E=Q\), but is
  excluded by any of Theorems 1--3 (e.g.~\(B(P,P)=a+b=0\neq 1\)).
\end{enumerate}

\emph{Proof.} Write
\(B(E,\rho)=a\,\mathrm{Tr}(E\rho)+b\,\mathrm{Tr}(E)\) with \(b=(1-a)/d\)
and \(\mathrm{Tr}\,\rho=1\). The two equations are
\[a\langle p|E|p\rangle+b\,\mathrm{Tr}(E)=1,\qquad a\langle q|E|q\rangle+b\,\mathrm{Tr}(E)=0.\]
Subtracting yields
\(a\bigl(\langle p|E|p\rangle-\langle q|E|q\rangle\bigr)=1\). For
\(0\leq E\leq I\) the difference of expectations on pure states lies in
\([-1,1]\). Hence \(|a|\geq 1\). Intersecting with the positivity
interval \(a\in[-1/(d-1),1]\):

\begin{itemize}
\tightlist
\item
  \(a=1\) is always admissible; take \(E=P\) when \(PQ=0\):
  \(B(P,P)=1\), \(B(P,Q)=b=0\).
\item
  \(a=-1/(d-1)\) requires the difference of expectations to equal
  \(1/a=-(d-1)\). This lies in \([-1,1]\) if and only if \(d-1\leq 1\),
  i.e.~\(d=2\). For \(d=2\), \(a=-1\), \(b=1\), and \(E=Q\) works:
  \(B(Q,P)=1\), \(B(Q,Q)=0\).
\end{itemize}

For \(d\geq 3\) the negative endpoint is therefore impossible under
distinguishability, leaving only \(a=1\). For \(d=2\) the residual
anti-Born point \(a=-1\) fails pure self-certainty (\(B(P,P)=0\neq 1\))
and fails the impossible-event condition in the Born direction
(\(B(P,Q)=1\neq 0\)). \(\blacksquare\)

\textbf{Remark.} Theorem 5 is operationally natural: existence of a
yes/no test that accepts one pure state with certainty and rejects an
orthogonal pure state with certainty. For qutrits and higher this single
requirement forces Born inside the bilinear class. The qubit retains a
single exotic anti-Born bilinear pairing, killed by the local selectors
of Theorems 1 and 3.

\textbf{Remark (failure of pure-state bits).} Equivalently: for
\(d\geq 3\) and \(a\in(-1/(d-1),1)\), \emph{no} pair of pure states is
perfectly distinguishable by any effect. The theory cannot encode a
classical bit in a pair of pure states with a yes/no test. This is a
sharp operational defect of the non-Born bilinear pairings, independent
of bit symmetry in the sense of Galley--Masanes {[}9{]} (which concerns
reversible maps between distinguishable pairs). Here the defect is more
primitive: the distinguishable pairs do not exist at all.

\subsubsection{\texorpdfstring{6.6 Failed candidate:
ancilla-independence (does \emph{not} select
Born)}{6.6 Failed candidate: ancilla-independence (does not select Born)}}\label{failed-candidate-ancilla-independence-does-not-select-born}

A natural candidate is ancilla-independence:
\(B_{dk}(E\otimes I_k,\rho\otimes\sigma)=B_d(E,\rho)\) for
\(\mathrm{Tr}\,\sigma=1\). With \(b_d=(1-a_d)/d\), the left-hand side
expands to
\[a_{dk}\,\mathrm{Tr}(E\rho)+\frac{1-a_{dk}}{d}\,\mathrm{Tr}(E).\]
Equality with the right-hand side for all \(E,\rho\) forces only
\(a_{dk}=a_d\), i.e.~stability of \(a\) under dimension scaling ---
\emph{not} \(a=1\). Ancilla-independence is therefore \emph{not} a Born
selector within this class. We record the failure to prevent a natural
mistake.

\subsubsection{6.7 Equivalence of the local
selectors}\label{equivalence-of-the-local-selectors}

\textbf{Proposition 3 (local selector equivalence).} Assume Lemma 2 and
normalization \(a+bd=1\) for \(d\geq 2\). The following are equivalent:

\begin{enumerate}
\def\labelenumi{\arabic{enumi}.}
\tightlist
\item
  there exist orthogonal rank-one projectors \(P,Q\) with \(B(Q,P)=0\)
  (hypothesis of Theorem 1);
\item
  \(B(P,P)=1\) for every rank-one projector \(P\) (hypothesis of Theorem
  3);
\item
  \(b=0\);
\item
  \(a=1\);
\item
  \(B(E,\rho)=\mathrm{Tr}(E\rho)\) for all effects and normalized
  states.
\end{enumerate}

\emph{Proof.} Under the classification, \(B(Q,P)=b\) for any orthogonal
rank-one pair and \(B(P,P)=a+b\) for any rank-one \(P\). Normalization
gives \(b=(1-a)/d\). Thus
(1)\(\Leftrightarrow b=0\Leftrightarrow a=1\Leftrightarrow a+b=1\Leftrightarrow\)(2)\(\Leftrightarrow\)(5).
\(\blacksquare\)

\textbf{Corollary.} Theorems 1 and 3 are two operational readings of the
\emph{same} algebraic condition \(b=0\). One is a single off-diagonal
vanishing (impossible event); the other is a single on-diagonal
certainty (pure self-prediction). They are not independent axioms once
bilinearity, unitary invariance, and normalization are granted.

The composition selector (Theorem 2), the dimension-independence
selector (Theorem 4), and the distinguishability selector (Theorem 5)
are \emph{independent routes} to the same conclusion (with the \(d=2\)
caveat of Theorem 5). Monoidality uses the tensor product structure of a
family \(\{B_d\}\); dimension-independence uses embeddings across
dimensions; distinguishability uses the existence of a perfect yes/no
test. Local selectors (Theorems 1 and 3) use only a single system and
are mutually equivalent.

\subsubsection{6.8 Status of the
selectors}\label{status-of-the-selectors}

The classification (Lemma 2) is standard representation theory. The
selectors are elementary: each is a one-line coefficient elimination
after the classification.

\begin{longtable}[]{@{}
  >{\raggedright\arraybackslash}p{(\linewidth - 8\tabcolsep) * \real{0.1667}}
  >{\raggedright\arraybackslash}p{(\linewidth - 8\tabcolsep) * \real{0.3500}}
  >{\raggedright\arraybackslash}p{(\linewidth - 8\tabcolsep) * \real{0.1333}}
  >{\raggedright\arraybackslash}p{(\linewidth - 8\tabcolsep) * \real{0.1667}}
  >{\raggedright\arraybackslash}p{(\linewidth - 8\tabcolsep) * \real{0.1833}}@{}}
\toprule\noalign{}
\begin{minipage}[b]{\linewidth}\raggedright
Selector
\end{minipage} & \begin{minipage}[b]{\linewidth}\raggedright
Operational content
\end{minipage} & \begin{minipage}[b]{\linewidth}\raggedright
Forces
\end{minipage} & \begin{minipage}[b]{\linewidth}\raggedright
Relation
\end{minipage} & \begin{minipage}[b]{\linewidth}\raggedright
Prior art
\end{minipage} \\
\midrule\noalign{}
\endhead
\bottomrule\noalign{}
\endlastfoot
Thm 1: one impossible event & \(B(Q,P)=0\) for one orthogonal pair &
\(b=0\) & \(\Leftrightarrow\) Thm 3 (Prop. 3) & apparently unrecorded in
this form \\
Thm 3: pure self-certainty & \(B(P,P)=1\) for pure \(P\) & \(a=1\) &
\(\Leftrightarrow\) Thm 1 (Prop. 3) & standard intuition; as
bilinear-class selector apparently unrecorded \\
Thm 2: rank-one monoidality & product probabilities factor & \(b=0\) (if
\(a\neq 0\)) & independent route & apparently unrecorded in this form \\
Thm 4: dim-independence & pure transitions ignore ambient \(d\) &
\(a_d=1\) & independent route & Hossenfelder {[}22{]} (stronger
setting) \\
Thm 5: perfect distinguishability & exists \(E\) with \(B(E,P)=1\),
\(B(E,Q)=0\) & \(a=1\) (\(d\geq 3\)); \(a\in\{1,-1\}\) (\(d=2\)) &
independent; \(d=2\) needs Thm 1/3 & apparently unrecorded in this
form \\
(failed) ancilla-independence & idle ancilla irrelevant & only
\(a_{dk}=a_d\) & does not select Born & --- \\
\end{longtable}

Based on targeted searches, Theorems 1--3 and 5 as Born selectors
\emph{within the invariant bilinear class} do not appear in the
literature we surveyed, but we have not performed a systematic review
and the results are elementary enough that they may exist as remarks or
exercises. Theorem 4 is a restricted restatement of the
dimension-independence idea of Hossenfelder {[}22{]}, who works directly
with continuous unitarily invariant transition probabilities on the
complex sphere (a stronger setup than the bilinear class). We claim
Theorems 1--3 and 5 as apparently-unrecorded corollaries, not as deep
new theorems.

Within the already-classified invariant bilinear class, the
orthogonal-support condition (Theorem 1) is substantially narrower than
full Gleason additivity: Gleason requires additivity over arbitrary
orthogonal decompositions, while Theorem 1 requires only vanishing on a
single orthogonal projector pair (extended to all such pairs by
invariance). We do not claim a global ordering of axiom strength between
the two frameworks, which operate under different background
assumptions.

\begin{center}\rule{0.5\linewidth}{0.5pt}\end{center}

\subsection{7. Limits and discussion}\label{limits-and-discussion}

\subsubsection{7.1 Countermodels}\label{countermodels}

The selectors exclude several theories, each by a distinct failure mode:

\begin{itemize}
\tightlist
\item
  \textbf{Classical probability:} no purification; the
  invariant-bilinear setup does not apply.
\item
  \textbf{Real-vector-space QM {[}11{]}:} fails local tomography; the
  Hermitian structure assumed here is absent.
\item
  \textbf{Quaternionic QM {[}5{]}:} fails local tomography;
  noncommutativity obstructs the tensor structure.
\item
  \textbf{Boxworld {[}13{]}:} the local gbit has four deterministic pure
  states and a square state space; it lacks purification in the strong
  sense. The Hermitian bilinear setup does not apply.
\item
  \textbf{GPTs without purification:} fail the existence of boundary
  lifts.
\end{itemize}

These exclusions are by the \emph{setup} (Hermitian structure assumed),
not by the selectors alone.

\subsubsection{7.2 What is genuinely here}\label{what-is-genuinely-here}

\begin{itemize}
\tightlist
\item
  Lemma 1: Gram map classification, continuous version (polynomial
  restriction is the \(k=1\) case of Procesi {[}2{]}).
\item
  Proposition 1: standard quotient geometry {[}3{]}.
\item
  Lemma 2: classification of invariant bilinear pairings via irreducible
  decomposition (standard representation theory).
\item
  \textbf{Theorems 1--5: selectors within the invariant bilinear class};
  Theorems 1 and 3 are equivalent (Proposition 3); Theorem 5 is the
  distinguishability selector (\(d\geq 3\) forces Born alone).
\item
  Theorem 4: dimension-independence step isolated from Hossenfelder
  {[}22{]}.
\item
  Explicit failure of ancilla-independence as a Born selector (Section
  6.6).
\item
  Honest supersession of Conjecture 1 as a Born-route by Nakahira
  {[}8{]} (Section 8); residual open Problem A (amplitude
  representations of OPTs).
\end{itemize}

\subsubsection{7.3 Relation to Gleason}\label{relation-to-gleason}

Gleason {[}14{]} derives the Born probability rule from noncontextual
additive measures on projectors (\(d\geq 3\)); Busch {[}15{]} extends to
POVMs including \(d=2\). Lemma 1 is complementary (state-formation map
vs probability rule). Theorem 1 is a \emph{different selector within a
different framework}: after the invariant bilinear classification, a
single vanishing condition selects Born. We do not claim Theorem 1
subsumes Gleason or is globally ``weaker''; only that within the
classified class, it is a narrower additional requirement.

\subsubsection{7.4 The circularity
boundary}\label{the-circularity-boundary}

Theorems 1--5 select Born \emph{given} the Hermitian structure of states
and effects, the unitary invariance, and the bilinear setup. They do not
derive these structures. The circularity is the same as Lemma 1: a
uniqueness result within an assumed geometry.

\subsubsection{7.5 Positioning relative to
Galley--Masanes}\label{positioning-relative-to-galleymasanes}

Galley and Masanes {[}9{]} keep pure states as rays and unitary
dynamics, and classify \emph{all} finite-parameter alternatives to the
measurement postulates. That landscape is large: non-Born theories
exist, and unrestricted ones violate bit symmetry. The present paper
studies only the \textbf{bilinear mixed-state extension} of outcome
probabilities --- the class containing the standard Born pairing
\(\mathrm{Tr}(E\rho)\) as a bilinear form on effects and states. Lemma 2
collapses that class to \(B_{a,b}\); the selectors finish the selection.
We do not claim to classify general alternatives to Born, nor to improve
on {[}9{]}. The contribution is a clean selector catalogue \emph{inside}
the bilinear slice, with explicit comparisons to Hossenfelder {[}22{]},
Gleason {[}14{]}, and Nakahira {[}8{]}.

\begin{center}\rule{0.5\linewidth}{0.5pt}\end{center}

\subsection{8. Confronting Conjecture 1 with
Nakahira}\label{confronting-conjecture-1-with-nakahira}

\textbf{Conjecture 1 (informal, original form).} In a finite-dimensional
causal operational theory admitting purification, with ancillary
stability and a suitable noncontextuality condition, the state-formation
map from amplitudes to states is forced to be congruence-covariant and
gauge-invariant, and hence (by Lemma 1) of the Gram form.

Nakahira {[}8{]} (arXiv:2605.23217, May 2026) proves a stronger
reconstruction theorem that must be confronted honestly.

\subsubsection{8.1 What Nakahira proves}\label{what-nakahira-proves}

In a finite-dimensional causal OPT, two postulates suffice:

\begin{enumerate}
\def\labelenumi{\arabic{enumi}.}
\tightlist
\item
  \textbf{Local equivalence:} channels are completely identified by
  local input-output statistics (equivalent to local discriminability).
\item
  \textbf{ES purification:} every state of system \(A\) is the marginal
  of a pure state on \(A\tilde A\) with \(\tilde A\) \emph{equivalent}
  to \(A\), and any two such purifications are related by a reversible
  channel on \(\tilde A\).
\end{enumerate}

\textbf{Theorem (Nakahira).} Under these two postulates the theory
\emph{is} standard finite-dimensional complex quantum theory: each
system has a finite-dimensional complex Hilbert space; states are
exactly all density matrices; measurements are exactly POVMs; composites
are tensor products; channels are exactly CPTP maps; the Born rule
holds; measurement no-restriction is \emph{derived}.

Neither postulate alone suffices: classical theory has local equivalence
but not ES purification; real quantum theory has ES purification but not
local equivalence (Nakahira Table I).

\subsubsection{8.2 What this does to Conjecture
1}\label{what-this-does-to-conjecture-1}

Nakahira does \textbf{not} construct the amplitude space
\(M_{d\times r}(\mathbb C)\) or the \(GL(d,\mathbb C)\times U(r)\)
bimodule. He constructs Hilbert spaces, density matrices, and the Born
pairing \(\mathrm{Tr}(E\rho)\) directly.

Consequences for our programme:

\begin{enumerate}
\def\labelenumi{\arabic{enumi}.}
\item
  \textbf{Born itself is no longer an open reconstruction target at the
  OPT level.} If Nakahira's theorem is accepted, the Born rule is
  already forced by local equivalence + ES purification. Our Theorems
  1--5 live one level \emph{down}: they select Born \emph{inside} the
  already Hermitian bilinear class. They remain valid mathematics, but
  they are not a competing derivation of Born from operational first
  principles.
\item
  \textbf{The amplitude/bimodule question is a representation question,
  not a reconstruction question.} Once Hilbert space is given, every
  density matrix of rank \(\leq r\) admits a Gram factorization
  \(\rho=WW^\dagger\) with \(W\in M_{d\times r}(\mathbb C)\), unique up
  to right \(U(r)\) (Proposition 1). Congruence covariance under
  invertible filters is the standard conjugation action. Lemma 1 then
  says: if one \emph{starts} from complex amplitudes and imposes gauge +
  congruence + positivity + continuity, the only state-formation map is
  the Gram map. That is a classification \emph{inside} the complex
  amplitude representation, not a derivation \emph{of} that
  representation.
\item
  \textbf{Revised open problem (replaces Conjecture 1).}\\
  \textbf{Problem A.} Characterize which OPTs admit an \emph{amplitude
  representation}: a functor assigning to each system a space
  \(M_{d\times r}\) with a \(GL(d)\times U(r)\) action whose quotient
  recovers the state space.\\
  \textbf{Problem B.} Among theories that already have Hilbert-space
  structure (e.g.~by Nakahira), is the Gram representation of mixed
  states forced by any additional operational principle beyond
  Stinespring/purification uniqueness, or is it merely a convenient
  coordinate system?\\
  Problem B is likely settled in the affirmative by standard dilation
  theory (Stinespring, Kraus): amplitudes are coordinates on dilations.
  The non-trivial direction is Problem A.
\item
  \textbf{What remains of our circularity boundary.} Lemma 1 and
  Theorems 1--5 never claimed to derive the Hermitian structure.
  Nakahira derives that structure from OPT postulates. The two results
  sit at different layers:

  \begin{itemize}
  \tightlist
  \item
    Nakahira: OPT postulates \(\to\) complex QM package (including
    Born);
  \item
    This paper: Hermitian bilinear invariant pairings \(\to\) Born via
    elementary selectors; and complex amplitudes + congruence \(\to\)
    Gram map.
  \end{itemize}
\end{enumerate}

\subsubsection{8.3 Honest status}\label{honest-status}

Conjecture 1 in its original form is \textbf{superseded} as a route to
Born: Nakahira already gets Born (and everything else) without
amplitudes. The residual mathematical content of Lemma 1 is a clean
classification of congruence-covariant continuous positive maps on
amplitude matrices --- valuable as linear algebra / invariant theory,
not as a foundational derivation of quantum probability. We keep Lemma 1
as scaffolding and reframe the open direction as Problem A above.

\begin{center}\rule{0.5\linewidth}{0.5pt}\end{center}

\subsection{9. Possible operational extension (not
formalized)}\label{possible-operational-extension-not-formalized}

We record, without formalizing, the operational intuition. Operational
non-uniqueness should be exhausted by reversible transformations on an
explicitly hidden interface. Two kinds of fibre arise: hidden-interface
fibres (from discarding a physical system \(E\); multiple realizations
related by reversible gauge on \(E\)) and observational fibres (from
restricting tests; trivial if the interface is complete). In quantum
theory, the first is essential uniqueness of purification and
Stinespring dilation; the second is local tomography. We do not
formalize a single principle covering both; in particular,
``completeness'' would need an independent definition (closure under
operations, realizability of dual-cone effects, universal property of
the observation functor) rather than being defined as injectivity, which
would be tautological. This is the formal content Conjecture 1 would
need.

\begin{center}\rule{0.5\linewidth}{0.5pt}\end{center}

\subsection{Acknowledgements}\label{acknowledgements}

The author used generative AI systems for conceptual exploration,
drafting, structural editing, and preliminary mathematical analysis. The
author reviewed and takes responsibility for the final manuscript.

\begin{center}\rule{0.5\linewidth}{0.5pt}\end{center}

\subsection{References}\label{references}

{[}1{]} H. Weyl, \emph{The Classical Groups: Their Invariants and
Representations}, Princeton University Press (1939).

{[}2{]} C. Procesi, ``Tensor fundamental theorems of invariant theory,''
arXiv:2011.10820 (2020).

{[}3{]} E. Massart and P.-A. Absil, ``Quotient geometry with simple
geodesics for the manifold of fixed-rank positive-semidefinite
matrices,'' \emph{SIAM J. Matrix Anal. Appl.} \textbf{41}(1), 171-198
(2020).

{[}4{]} G. Chiribella, G. M. D'Ariano, and P. Perinotti, ``Informational
derivation of quantum theory,'' \emph{Phys. Rev.~A} \textbf{84}, 012311
(2011). arXiv:1011.6451.

{[}5{]} H. Barnum and A. Wilce, ``Local tomography and the Jordan
structure of quantum theory,'' \emph{Found. Phys.} \textbf{44}, 946-979
(2014). arXiv:1202.4513.

{[}6{]} S. Tull, ``A categorical reconstruction of quantum theory,''
\emph{Log. Methods Comput. Sci.} \textbf{16}, 1 (2020).
arXiv:1804.02265.

{[}7{]} J. H. Selby, C. M. Scandolo, and B. Coecke, ``Reconstructing
quantum theory from diagrammatic postulates,'' arXiv:1802.00367 (2018).

{[}8{]} K. Nakahira, ``Two operational principles single out quantum
theory,'' arXiv:2605.23217 (2026, preprint).

{[}9{]} T. D. Galley and L. Masanes, ``Classification of all
alternatives to the Born rule in terms of informational properties,''
\emph{Phys. Rev.~Lett.} \textbf{120}, 250403 (2018). arXiv:1610.04859.

{[}10{]} J. Zhang, ``Summing to Uncertainty: On the Necessity of
Additivity in Deriving the Born Rule,'' arXiv:2603.06211 (2026,
preprint).

{[}11{]} L. Hardy, ``Limited holism and real-vector-space quantum
theory,'' \emph{Found. Phys.} \textbf{42}, 454-471 (2012).
arXiv:1005.4870.

{[}12{]} J. Bang, K. Cho, and K. Baek, ``Hidden complex structure in
quotient-space real quantum mechanics,'' arXiv:2607.05865 (2026,
preprint).

{[}13{]} J. Barrett, ``Information processing in generalized
probabilistic theories,'' \emph{Phys. Rev.~A} \textbf{75}, 032304
(2007). arXiv:quant-ph/0508211.

{[}14{]} A. Gleason, ``Measures on the closed subspaces of a Hilbert
space,'' \emph{J. Math. Mech.} \textbf{6}, 885-893 (1957).

{[}15{]} P. Busch, ``Quantum states and generalized observables: A
simple proof of Gleason's theorem,'' \emph{Phys. Rev.~Lett.}
\textbf{91}, 120403 (2003). arXiv:quant-ph/9909073.

{[}16{]} S. S. Holland Jr., ``Orthomodularity in infinite dimensions; a
theorem of M. Solèr,'' arXiv:math/9504224 (1995).

{[}17{]} W. M. Stuckey, M. Silberstein, and T. McDevitt, ``Completing
the Quantum Reconstruction Program via the Relativity Principle,''
arXiv:2404.13064 (2024).

{[}18{]} J. van de Wetering, ``Self-duality and Jordan structure of
quantum theory follow from homogeneity and pure transitivity,''
arXiv:2306.00362 (2023).

{[}19{]} S. Tull, ``Categorical operational physics,'' arXiv:1902.00343
(2019).

{[}20{]} O. Cunningham and C. Heunen, ``Axiomatizing complete
positivity,'' \emph{Electronic Proceedings in Theoretical Computer
Science} \textbf{195}, 148-157 (2015). arXiv:1506.02931.

{[}22{]} S. Hossenfelder, ``A derivation of Born's rule from symmetry,''
\emph{Ann. Phys.} \textbf{425}, 168394 (2021). arXiv:2006.14175.

{[}23{]} A. Lax, ``Fixed-PVM Born Rule Uniqueness from Fisher
Non-Expansion and Operational Calibration,'' arXiv:2604.27339 (2026,
preprint).

\begin{center}\rule{0.5\linewidth}{0.5pt}\end{center}

\emph{End of manuscript.}

\end{document}
